% GENERATED FILE. Edit canonical lessons, then run npm run generate:books.
% canonical-source-hash: fnv1a64:96bb80efc3174ee7
% canonical-lessons: RU-C214-otsyuda, RU-C214-nazad, RU-C214-vperyod, RU-C214-vverkh, RU-C214-vniz

\chapter{From here, Back, Forward, Up, Down}
\label{ch:ru-from-here-back-forward-up-down}
\hlchaptermodality{214}

\begin{chapteropening}
\textbf{By the end of this chapter:} \emph{I can say from here, back, forward, up and down.}

Complete the last lesson of chapter 214: I can say from here, back, forward, up and down.
\end{chapteropening}

\section[\texorpdfstring{\ru{отсюда} (otsyúda)}{отсюда (otsyúda)}]{\ru{отсюда} (otsyúda) — from here}

\label{lesson:RU-C214-otsyuda}



\begin{quote}
Before the new one: say the Russian for a flat, then the Russian for a floor.
\end{quote}

\subsection*{You'll want to know: \ru{отсюда}}
\textbf{\ru{отсюда}} (\emph{otsyúda}) — \textquotedblleft{}from here\textquotedblright{}.

\begin{grammarlens}[title={What you've built}]
One more word for this chapter.
\end{grammarlens}

\subsection*{Your turn}
\begin{itemize}
\raggedright
  \item \emph{Say it:} \emph{otsyúda}
  \item \emph{Say it:} \emph{otsyúda}, once more
  \item \emph{Say it:} say \emph{otsyúda}
  \item \emph{Recall:} say the Russian for a flat, then the Russian for a floor, then say \emph{otsyúda} again
\end{itemize}

\begin{tcolorbox}[breakable,colback=teal!4,colframe=teal!35!black,title={Before you move on}]
What does \ru{отсюда} mean? (\textquotedblleft{}From here\textquotedblright{}.) Say it once more.
\end{tcolorbox}
\section[\texorpdfstring{\ru{назад} (nazád)}{назад (nazád)}]{\ru{назад} (nazád) — back, backwards}

\label{lesson:RU-C214-nazad}



\begin{quote}
Before the new one: say the Russian for a floor, then the Russian for from here.
\end{quote}

\subsection*{You'll want to know: \ru{назад}}
\textbf{\ru{назад}} (\emph{nazád}) — \textquotedblleft{}back, backwards\textquotedblright{}.

\begin{grammarlens}[title={What you've built}]
One more word for this chapter.
\end{grammarlens}

\subsection*{Your turn}
\begin{itemize}
\raggedright
  \item \emph{Say it:} \emph{nazád}
  \item \emph{Say it:} \emph{nazád}, once more
  \item \emph{Say it:} say \emph{nazád}
  \item \emph{Recall:} say the Russian for a floor, then the Russian for from here, then say \emph{nazád} again
\end{itemize}

\begin{tcolorbox}[breakable,colback=teal!4,colframe=teal!35!black,title={Before you move on}]
What does \ru{назад} mean? (\textquotedblleft{}Back, backwards\textquotedblright{}.) Say it once more.
\end{tcolorbox}
\section[\texorpdfstring{\ru{вперёд} (vperyód)}{вперёд (vperyód)}]{\ru{вперёд} (vperyód) — forward}

\label{lesson:RU-C214-vperyod}



\begin{quote}
Before the new one: say the Russian for from here, then the Russian for back.
\end{quote}

\subsection*{You'll want to know: \ru{вперёд}}
\textbf{\ru{вперёд}} (\emph{vperyód}) — \textquotedblleft{}forward\textquotedblright{}.

\begin{grammarlens}[title={What you've built}]
One more word for this chapter.
\end{grammarlens}

\subsection*{Your turn}
\begin{itemize}
\raggedright
  \item \emph{Say it:} \emph{vperyód}
  \item \emph{Say it:} \emph{vperyód}, once more
  \item \emph{Say it:} say \emph{vperyód}
  \item \emph{Recall:} say the Russian for from here, then the Russian for back, then say \emph{vperyód} again
\end{itemize}

\begin{tcolorbox}[breakable,colback=teal!4,colframe=teal!35!black,title={Before you move on}]
What does \ru{вперёд} mean? (\textquotedblleft{}Forward\textquotedblright{}.) Say it once more.
\end{tcolorbox}
\section[\texorpdfstring{\ru{вверх} (vverkh)}{вверх (vverkh)}]{\ru{вверх} (vverkh) — up, upwards}

\label{lesson:RU-C214-vverkh}



\begin{quote}
Before the new one: say the Russian for back, then the Russian for forward.
\end{quote}

\subsection*{You'll want to know: \ru{вверх}}
\textbf{\ru{вверх}} (\emph{vverkh}) — \textquotedblleft{}up, upwards\textquotedblright{}.

\begin{grammarlens}[title={What you've built}]
One more word for this chapter.
\end{grammarlens}

\subsection*{Your turn}
\begin{itemize}
\raggedright
  \item \emph{Say it:} \emph{vverkh}
  \item \emph{Say it:} \emph{vverkh}, once more
  \item \emph{Say it:} say \emph{vverkh}
  \item \emph{Recall:} say the Russian for back, then the Russian for forward, then say \emph{vverkh} again
\end{itemize}

\begin{tcolorbox}[breakable,colback=teal!4,colframe=teal!35!black,title={Before you move on}]
What does \ru{вверх} mean? (\textquotedblleft{}Up, upwards\textquotedblright{}.) Say it once more.
\end{tcolorbox}
\section[\texorpdfstring{\ru{вниз} (vniz)}{вниз (vniz)}]{\ru{вниз} (vniz) — down, downwards}

\label{lesson:RU-C214-vniz}



\begin{quote}
Before the new one: say the Russian for forward, then the Russian for up.
\end{quote}

\subsection*{You'll want to know: \ru{вниз}}
\textbf{\ru{вниз}} (\emph{vniz}) — \textquotedblleft{}down, downwards\textquotedblright{}.

\begin{grammarlens}[title={What you've built}]
One more word for this chapter.
\end{grammarlens}

\subsection*{Your turn}
\begin{itemize}
\raggedright
  \item \emph{Say it:} \emph{vniz}
  \item \emph{Say it:} \emph{vniz}, once more
  \item \emph{Say it:} say \emph{vniz}
  \item \emph{Recall:} say the Russian for forward, then the Russian for up, then say \emph{vniz} again
\end{itemize}

\begin{tcolorbox}[breakable,colback=teal!4,colframe=teal!35!black,title={Before you move on}]
What does \ru{вниз} mean? (\textquotedblleft{}Down, downwards\textquotedblright{}.) Say it once more.
\end{tcolorbox}
